\documentclass[11pt]{article}
\usepackage[T1]{fontenc}
\usepackage{lmodern}
\usepackage[margin=1in]{geometry}
\usepackage{amsmath,amssymb,amsthm,booktabs}
\usepackage{microtype}
\usepackage[hidelinks]{hyperref}
\hypersetup{pdftitle={A coefficient obstruction for antiprism initial flag determinants},
pdfauthor={Xiaojian Zeng}}
\newtheorem{theorem}{Theorem}
\newtheorem{lemma}[theorem]{Lemma}
\newcommand{\Q}{\mathbb{Q}}
\newcommand{\tr}{\operatorname{tr}}
\newcommand{\diag}{\operatorname{diag}}
\newcommand{\fall}[2]{(#1)_{\underline{#2}}}
\newcommand{\coef}[1]{\left[#1\right]}
\title{A coefficient obstruction for antiprism initial flag determinants}
\author{Xiaojian Zeng\\[3pt]
\small PBC School of Finance, Tsinghua University, 100190\\
\small\texttt{xiaojianzeng@foxmail.com}}
\date{September 14, 2026}
\begin{document}
\maketitle

\begin{abstract}
The total nonnegativity question for antiprism $h$-vector transformation
matrices remains unresolved. We study a specific stronger certificate:
coefficient nonnegativity of a formal determinant whose specialization at
nonnegative parameters gives consecutive initial flags. We compute a family of
coefficients in its next-highest degree and obtain an explicit quartic
formula that is negative at every order at least $76$. An exact negative
coefficient already occurs at order $22$, while a second family is positive
at every order. These identities rule out the proposed certificate in the
specified variables. They do not provide a negative minor of an antiprism
matrix or an obstruction to other proofs of total nonnegativity.
\end{abstract}

\section{The matrix question and the auxiliary determinant}

Write $S(u,v)$ for a Stirling number of the second kind, with $S(0,0)=1$.
Binomial coefficients and Stirling numbers with invalid indices are zero.
For $d\geq1$ and $0\leq i,j\leq d$, define
\begin{equation}\label{eq:H}
 H_d(i,j)=\sum_{r=0}^d\sum_{s=0}^i
 \binom{j}{s}\binom{d-j}{d-r-s}\binom{r}{i-s}
 \sum_{\ell\geq0}(\ell!)^2S(i-s,\ell)S(r-i+s,\ell).
\end{equation}
The sums are finite. This is the matrix taking the $h$-vector of a
$(d-1)$-dimensional simplicial complex to that of its antiprism subdivision.
Athanasiadis, Brunink and Juhnke-Kubitzke give the transformation and its
combinatorial interpretation in Section~5.3 of \cite{ABJ2022}; see also
their Proposition~5.7(b) for the local Stirling formula.

A matrix is \emph{totally nonnegative} if every minor, with its row and
column indices in increasing order, is nonnegative. The subdivision
literature cited here calls this weak inequality total positivity.
Mu and Welker report positive antiprism experiments in small dimensions
\cite[Section~2, final paragraph]{MW2026}.
Liu and Mao prove total nonnegativity for barycentric and interval
subdivision matrices and raise the corresponding antiprism question
\cite[Remark~5.8 in the PDF]{LM2026}. Their HTML version numbers that
remark~5.20. The original question is whether every $H_d$ is totally
nonnegative, for every $d\geq1$; this note leaves it unresolved.

We consider the flags with rows $1,\ldots,k$ and columns
$j-k+1,\ldots,j$, where $d\geq j\geq k\geq1$. The following formal
polynomial arises by replacing the exponential factors in these flags by
independent partial sums. All variables $a,n,z_1,\ldots,z_k$ below are
algebraically independent. Put
\begin{equation}\label{eq:T}
 T_{r,\ell}(x,y)=\coef{t^r}e^{-ty}
       (1-e^{-t(x+y)})^\ell,\qquad
 w_\ell=\sum_{h=1}^{\ell}z_h,\qquad q_\ell=\frac{\ell}{\ell+1}.
\end{equation}
Here $T_{r,\ell}$ is homogeneous of degree $r$ and is zero if $\ell>r$.
With $\fall{v}{p}=\prod_{h=0}^{p-1}(v-h)$ and $\fall{v}{0}=1$, set
\begin{align}
 M_{r,s}&=\sum_{\ell=1}^r\ell^k w_\ell q_\ell^s
 \sum_{p+q=r}\coef{x^py^q}T_{r,\ell}(x,y)
 \frac{\fall{a+k+s}{p}\fall{n-s}{q}}{\ell^p(\ell+1)^q},
 \label{eq:M}\\
 R_k&=\det(M_{r,s})_{\substack{1\leq r\leq k\\1-k\leq s\leq0}},
 \qquad K=\frac{k(k+1)}2,\qquad
 \lambda_k=\frac{\prod_{r=1}^k r!}{\prod_{\ell=1}^k\ell^\ell}>0.
 \label{eq:R}
\end{align}
Both index orders in this determinant are increasing, and $p,q$ in
\eqref{eq:M} are nonnegative integers. In particular,
$R_k\in\Q[a,n,z_1,\ldots,z_k]$.

\begin{theorem}\label{thm:main}
For every integer $k\geq5$,
\begin{equation}\label{eq:quartic}
 \coef{a^{K-1}n^0z_1^5z_6\cdots z_k}R_k
 =-\frac{\lambda_k}{20736}
  \bigl(125k^4-9555k^3+17736k^2-346478k+858264\bigr).
\end{equation}
The final product is $1$ when $k=5$. This coefficient is negative for
every integer $k\geq76$. In addition,
\begin{equation}\label{eq:22}
 \coef{a^{252}n^0z_1^9z_{10}\cdots z_{22}}R_{22}
 =-\lambda_{22}\frac{4429926944297146258237}{14051921264640000000}<0.
\end{equation}
For every integer $k\geq1$,
\begin{equation}\label{eq:positive}
 \coef{a^{K-1}n^0z_1z_2\cdots z_k}R_k
 =\lambda_k\frac{k(k+1)(2k+1)}6>0.
\end{equation}
Thus the homogeneous component of total $(a,n)$-degree $K-1$ has
coefficients of both signs for every $k\geq76$, and also for $k=22$.
\end{theorem}

The assertion that all coefficients of this particular $R_k$ are
nonnegative was a working hypothesis introduced in the present
investigation, motivated by exact low-order expansions. It is not a
conjecture attributed to any of the cited authors. The theorem identifies
an obstruction to that specific certificate. Our comparison with the
primary subdivision papers and bounded searches through September~14,
2026 found no equivalent coefficient formula; this is a limited literature
assessment, not an exhaustive priority claim.

To verify the connection with \eqref{eq:H}, extend the notation by
$H_0(0,0)=1$ and form the formal generating series
\begin{equation}\label{eq:G}
 G(x,y,t)=\sum_{j,n\geq0}\sum_{r=0}^{j+n}
 H_{j+n}(r,j)t^r\frac{x^jy^n}{j!n!}
 =\frac{e^{tx+y}}{1-(e^{t(x+y)}-1)(e^{x+y}-1)}.
\end{equation}
Indeed, in \eqref{eq:H} put $p=i-s$ and $q=r-i+s$. Summing the three
binomial factors against $x^jy^n/(j!n!)$ gives
$x^s e^y(x+y)^{p+q}/(s!p!q!)$.
Summing $s$, then using
$\sum_{u\geq0}S(u,\ell)v^u/u!=(e^v-1)^\ell/\ell!$
for each of $p,q$, proves \eqref{eq:G}.
Rewriting its denominator gives
\[
 \coef{t^r}G=\sum_{\ell=0}^r
 T_{r,\ell}(x,y)e^{\ell x+(\ell+1)y}.
\]
The $\ell=0$ term has no $x$ dependence. For $j+s\geq1$, exponential
coefficient extraction, multiplied by $(j+s)!(d-j-s)!$, therefore gives
$M_{r,s}=H_d(r,j+s)$ under the specialization
\begin{align}
 a&=j-k,\qquad n=d-j,\qquad z_1=2^n,\label{eq:special}\\
 z_\ell&=\ell^a(\ell+1)^n-(\ell-1)^a\ell^n\quad(2\leq\ell\leq k).
 \notag
\end{align}
All these parameters are nonnegative, and the partial sums telescope to
$w_\ell=\ell^a(\ell+1)^n$. All series manipulations are formal and locally
finite. Coefficient nonnegativity of $R_k$ would consequently certify
nonnegativity of these flags. A negative formal coefficient need not give
a negative value at \eqref{eq:special}.

\section{The next coefficient in the row-degree expansion}

Since every coefficient in Theorem~\ref{thm:main} has $n$-exponent zero,
we may set $n=0$. Define constant lower triangular matrices $A,E,F$ by
\begin{align}
 A_{r,\ell}&=\frac{\coef{x^r}T_{r,\ell}}{\ell^r},&
 E_{r,\ell}&=rA_{r,\ell}
 -\frac{\coef{x^{r-1}y}T_{r,\ell}}{\ell^{r-1}(\ell+1)},&
 F&=A^{-1}E.\label{eq:AE}
\end{align}
The diagonal entries of $A$ are $\ell^{-\ell}$, so $A$ is invertible.
Let
\[
 B=\diag(\ell^k w_\ell)_{\ell=1}^k,\quad
 V=(q_\ell^s)_{\ell,s},\quad S=\diag(1-k,\ldots,0),\quad
 D=\diag\left(kr-\frac{r(r-1)}2\right)_{r=1}^k.
\]
In row $r$ of $M|_{n=0}$, the coefficient of $a^r$ is the corresponding
row of $U=ABV$, and that of $a^{r-1}$ is the row of
\begin{equation}\label{eq:N}
 N=DABV+EBVS.
\end{equation}
For the latter assertion, the $q=0$ term of \eqref{eq:M} contributes
$r(k+s)-r(r-1)/2$ from the falling factorial. The $q=1$ term contributes
$-s$; all $q\geq2$ terms have $a$-degree at most $r-2$.
No other term contributes to degree $K-1$ in the determinant.

Factoring $q_\ell^{1-k}$ from each row of $V$ reduces its determinant to
an ordinary Vandermonde determinant. Since
$q_v-q_u=(v-u)/((u+1)(v+1))$ for $u<v$, this gives
\begin{equation}\label{eq:detU}
 \det V=\frac{\prod_{p<q}(q-p)}{\prod_{\ell=1}^k\ell^{k-1}},
 \qquad \det U=\lambda_k\prod_{\ell=1}^k w_\ell.
\end{equation}
In the field $\Q(z_1,\ldots,z_k)$, determinant multilinearity now yields
\[
 \coef{a^{K-1}n^0}R_k=\det U\,\tr(U^{-1}N).
\]
Put $W=VSV^{-1}$. Cyclicity of trace and lower triangularity of $F$ give
\begin{equation}\label{eq:trace}
 \tr(U^{-1}N)=\tr D+
 \sum_{1\leq\ell\leq i\leq k}
 F_{i,\ell}W_{\ell,i}\left(\frac\ell i\right)^k\frac{w_\ell}{w_i}.
\end{equation}
Multiplication by $\prod w_\ell$ clears every denominator involving $z$,
so the resulting identity is polynomial and also holds when some $w_i=0$.

We next evaluate the needed entries of $W$. Write
$\mathcal H_v=\sum_{h=1}^v1/h$, with $\mathcal H_0=0$. Then
\begin{align}
 W_{\ell,\ell}
 &= (\ell+1)\bigl(\ell(\mathcal H_{\ell-1}-\mathcal H_{k-\ell})-(k-1)\bigr),
 \label{eq:Wdiag}\\
 \left(\frac\ell i\right)^k W_{\ell,i}
 &=(-1)^{i-\ell+1}\frac{\ell(\ell+1)}{i-\ell}
 \frac{\binom{k-\ell}{i-\ell}}{\binom{i}{i-\ell}}\quad(\ell<i).
 \label{eq:Woff}
\end{align}
Here is a direct interpolation derivation, including the sign.
Let $P(q)=\prod_{h=1}^k(q-q_h)$ and
$L_i(q)=P(q)/((q-q_i)P'(q_i))$.
These are the usual Lagrange cardinal polynomials; differentiation of this
representation is the standard construction of an interpolation
differentiation matrix \cite[Sections~2 and~9.3]{BT2004}.
The reciprocal derivatives $1/P'(q_i)$ also occur as the values of the
canonical inverse of $P'$ modulo $P$ in Lemma~2.1 of
Zeng \cite{Zeng2026}. We use this elementary
weight identity, without a complexity assertion.

For the Laurent basis $q^{1-k},\ldots,1$, the cardinal functions are
$C_i(q)=(q_i/q)^{k-1}L_i(q)$. Hence
\[
 W_{\ell,i}=q_\ell C_i'(q_\ell),\qquad
 W_{\ell,i}=q_\ell(q_i/q_\ell)^{k-1}L_i'(q_\ell)\quad(\ell\ne i).
\]
Inserting $q_\ell-q_h=(\ell-h)/((\ell+1)(h+1))$ gives
\eqref{eq:Woff}. On the diagonal, the same calculation gives
$W_{\ell,\ell}=1-k+q_\ell\sum_{h\ne\ell}(q_\ell-q_h)^{-1}$,
which simplifies to \eqref{eq:Wdiag}.
Also $F_{\ell,\ell}=\ell/(\ell+1)$. Thus the diagonal contribution
to \eqref{eq:trace} is
\begin{align}
 \tr D+\sum_{\ell=1}^kF_{\ell,\ell}W_{\ell,\ell}
 &=\sum_{\ell=1}^k\left\{\frac{\ell(3-\ell)}2+
 \ell^2(\mathcal H_{\ell-1}-\mathcal H_{k-\ell})\right\}\notag\\
 &=\frac{k(k+1)(2k+1)}6.\label{eq:diagsum}
\end{align}
To check the cancellation, pair the harmonic terms with indices $p<q$:
their sum is $\sum_{p<q}(q^2-p^2)/(q-p)=\sum_{p<q}(p+q)$.

\begin{lemma}\label{lem:family}
For $1\leq m\leq k$, put $Z_{k,m}=z_1^m z_{m+1}\cdots z_k$.
Then
\begin{equation}\label{eq:family}
 \frac{\coef{a^{K-1}n^0Z_{k,m}}R_k}{\lambda_k}
 =\frac{k(k+1)(2k+1)}6+
 \sum_{1\leq\ell<i\leq m}
 (-1)^{i-\ell+1}\frac{\ell(\ell+1)}{i-\ell}
 \frac{\binom{k-\ell}{i-\ell}}{\binom{i}{i-\ell}}F_{i,\ell}.
\end{equation}
\end{lemma}
\begin{proof}
The coefficient of $Z_{k,m}$ in $\prod w_h$ is $1$. For $\ell<i$, its
coefficient in $(\prod w_h)w_\ell/w_i$ is $1$ if $i\leq m$, and $0$
otherwise. Indeed, $z_k$ must come from $w_k$; continue downwards through
$z_{m+1}$. A missing factor $w_i$ in this range makes the assignment
impossible. Otherwise all remaining factors must supply $z_1$, uniquely.
Apply this observation to \eqref{eq:trace}, using
\eqref{eq:detU}--\eqref{eq:diagsum}.
\end{proof}

\section{Evaluation of the negative coefficients}

The entries needed in \eqref{eq:AE} admit the explicit formulas
\begin{align}
 A_{r,\ell}&=\frac{(-1)^{r-\ell}\ell!S(r,\ell)}{r!\ell^r},&
 E_{r,\ell}&=\frac{(-1)^{r-\ell}(\ell-1)!S(r-1,\ell-1)}
 {(\ell+1)(r-1)!\ell^{r-1}}.\label{eq:stirlingAE}
\end{align}
The first follows from \eqref{eq:T}. For the second, apply
$\partial_a-\partial_n$ to
$e^{-tn/(\ell+1)}(1-e^{-t(a/\ell+n/(\ell+1))})^\ell$.
The result is
\[
 \frac{t}{\ell+1}e^{-tn/(\ell+1)}
 (1-e^{-t(a/\ell+n/(\ell+1))})^{\ell-1}.
\]
Take the $t^r$ coefficient at $(a,n)=(1,0)$.
Forward substitution in $AF=E$ gives the first five rows
\begin{equation}\label{eq:Ftable}
 \begin{pmatrix}
 1/2&0&0&0&0\\
 1&2/3&0&0&0\\
 9/8&9/8&3/4&0&0\\
 4/3&8/9&32/27&4/5&0\\
 625/432&3125/3456&3125/3888&625/512&5/6
 \end{pmatrix}.
\end{equation}
Lower triangularity makes these entries independent of $k$ once $k\geq5$.

\begin{proof}[Proof of Theorem~\ref{thm:main}]
For $m=5$, substitute \eqref{eq:Ftable} into \eqref{eq:family} and expand
the binomial polynomials. The result is \eqref{eq:quartic}. This is a
polynomial identity of degree at most four, obtained from ten explicit
rational terms. If $Q(k)$ is its quartic numerator, then
\begin{equation}\label{eq:shift}
 Q(76+t)=125t^4+28445t^3+2171196t^2+56268354t+52825392.
\end{equation}
Every coefficient is positive, so $Q(k)>0$ for real $k\geq76$.
This proves the uniform negative sign.

For completeness, the larger arithmetic in \eqref{eq:22} is specified
by a short recurrence and table. Compute $S(r,\ell)$ using
$S(r,\ell)=S(r-1,\ell-1)+\ell S(r-1,\ell)$ through $r=9$, and compute
\begin{equation}\label{eq:forward}
 F_{r,\ell}=\frac{E_{r,\ell}-\sum_{h=\ell}^{r-1}A_{r,h}F_{h,\ell}}
 {A_{r,r}}\quad(1\leq\ell\leq r\leq9)
\end{equation}
from \eqref{eq:stirlingAE}. For $k=22$, group the summands in
\eqref{eq:family} by $i$, calling the resulting sum $c_i$.
The exact values are in Table~\ref{tab:22}.
\begin{table}[htbp]
\centering
\caption{Off-diagonal contributions $c_i$ for $k=22$, $m=9$.}
\label{tab:22}
\begin{tabular}{cr}
\toprule
$i$ & $c_i$\\
\midrule
2 & $21$\\
3 & $-135/4$\\
4 & $836/3$\\
5 & $-41875/64$\\
6 & $34868547/20000$\\
7 & $-1924165512913/583200000$\\
8 & $315237249593344/53603825625$\\
9 & $-155107959187079458053/19275612160000000$\\
\bottomrule
\end{tabular}
\end{table}
Since $22\cdot23\cdot45/6=3795$, exact summation gives
\[
 3795+\sum_{i=2}^9 c_i
 =-\frac{4429926944297146258237}{14051921264640000000}.
\]
Lemma~\ref{lem:family}, with $K-1=252$, proves \eqref{eq:22}.
Finally, taking $m=1$ in that lemma leaves an empty off-diagonal sum and
proves \eqref{eq:positive}, including $k=1$.
\end{proof}

\section{Exact verification and limits}

The companion Python script \texttt{verify.py} computes the displayed
tables and checks the quartic identity with integer and rational
arithmetic. It also checks the order-$22$ coefficient directly from
\eqref{eq:M}, without using $F$, $W$, or \eqref{eq:family} for that
extraction. Here is the completeness argument for the second calculation.
In a determinant term contributing $Z_{k,m}$, row $k$ must supply $z_k$,
then row $k-1$ must supply $z_{k-1}$, down through row $m+1$.
Consequently keep only $\ell=r$ in rows $r>m$ of \eqref{eq:M}, and all
$\ell\leq r$ in the first $m$ rows, removing the chosen $z$ factors.
The determinant of the resulting row vectors is exactly
$\coef{Z_{k,m}}R_k$.

At $n=0$, its two leading row coefficients use only
$\coef{x^r}T_{r,\ell}$ and $\coef{x^{r-1}y}T_{r,\ell}$. These may be
computed directly from
\begin{equation}\label{eq:powersum}
 \coef{x^py^{r-p}}T_{r,\ell}
 =\frac{(-1)^r}{p!(r-p)!}
 \sum_{h=0}^{\ell}(-1)^h\binom{\ell}{h}h^p(h+1)^{r-p},
\end{equation}
with $0^0=1$. Clear the rational denominators in each pair of row
coefficients. The coefficient of $a^{K-1}$ is the sum of the integer
determinants obtained by replacing one leading row by its next row,
divided by the product of the clearing denominators. The implementation
uses fraction-free elimination of Bareiss \cite[Section~I]{Bareiss1968},
checking every division for zero remainder. This gives the same value
in \eqref{eq:22}.

The exploratory search covered only
$\coef{a^{K-1}n^0Z_{k,m}}R_k$ for $2\leq k\leq32$ and
$1\leq m\leq k$. Its first negative order was $22$ within that family
and range. No minimality among all coefficients is asserted. The earlier
complete rational expansions at orders five and six have, respectively,
$5712$ and $33396$ nonzero coefficients, all positive; these finite
observations motivated the certificate and are not premises of the proof.

The obstruction applies to coefficient nonnegativity in the unchanged
variables of \eqref{eq:M}. The exponential relations in
\eqref{eq:special} may still permit nonnegative flag values. Moreover,
these flags cover only selected rows and consecutive positive columns;
they do not by themselves settle all minors of $H_d$. Alternative
parametrizations and other total-nonnegativity arguments remain possible.

\paragraph{Computational and AI assistance.}
OpenAI Codex assisted with the symbolic derivation, exact-arithmetic
programs, literature checks, and manuscript preparation. A separate model
context checked the proof and recomputed the finite witness by exact
elimination over $\Q[\varepsilon]/(\varepsilon^2)$. That check is not
external peer review. No floating-point calculation or proof-assistant
formalization is used to establish the theorem. The companion script
requires only the Python standard library; the separate verification
also used SymPy~1.14.0.

\bibliographystyle{plain}
\bibliography{references}
\end{document}
